\begin{definition} \textit{Конфигурационным графом} для машины $M$ и входа $x$ называется граф, вершинами которого являются все конфигурации для входа $x$, а рёбрами --- возможность перехода из одной конфигурации в другую за один шаг работы $M$ на входе $x$. Для определённости будем считать, что в графе есть ровно одна принимающая конфигурация (например, добавим её искусственно и пустим в неё рёбра из всех конфигураций с принимающим состоянием).
\end{definition}

\begin{note} В таком случае (недетерминированная) машина принимает слово тогда и только тогда, когда в графе есть путь из начальной вершины в принимающую.
\end{note}

\begin{theorem}[Сэвич] Если $s(n) \ge \log n$, то $\nsp{s(n)} \subset \dsp{s(n)^2}$.
\end{theorem}

\begin{proof} Рассмотрим конфигурационный граф. Нам нужно понять, есть ли пусть между двучмя фиксированными вершинами $u$ и $v$ в графе экспоненциального размера $2^n$. Заметим, что если путь есть, то есть и путь длины не больше числа вершин в графе. 

Обозначим через $R(x, y, k)$ предикат существования пути не длиннее $2^k$ из вершины $x$ в вершину $y$. В таком случае имеем следующее рекурсивное соотношение: $R(x, y, k)$ верно тогда и только тогда, когда для некоторой вершины $z$ выполнены одновременно $R(x, z, k-1)$ и $R(z, y, k-1)$. Таким образом, получается следующий рекурсивный алгоритм: перебираем все $z$, для каждого вычисляем истинность $R(x, z, k-1)$, запоминаем результат и на той же памяти вычисляем $R(z, y, k-1)$. Далее возвращаем конъюнкцию двух результатов.

Базой рекурсии будет случай $k=0$: в этом случае нужно проверить, что либо $x$ и $y$ совпадают, либо из $x$ в $y$ можно попасть за один шаг машины Тьюринга (это требует совсем мало дополнительной памяти).

На каждом шаге рекурсии требуется запомнить только текущую вершину $z$: это требует $O(s(n))$ памяти (здесь мы и используем условие $s(n) \geq \log n$: только в таком случае число вершин в конфигурационном графе будет равно $O(n2^s)$, то есть $2^{O(s(n))}$). Глубина рекурсии также составит $O(s(n))$. Итого, общая использованная память равна $O(s(n)^2)$, ч.т.д.
\end{proof}

\begin{corollary}
$\psp = \npsp$.
\end{corollary}

\begin{proof} Применим теорему Сэвича и заметим, что квадрат полинома также является полиномом.
\end{proof}